% P05.2 local rank-one graph orthogonality proof, 25 September 2026.

\begin{proof}[\Cref{thm:frontier-polyhedral-local-rank-one-graph}]
Fix $u\in\mathrm{PA}_{\mathbf R,b}(\mathbf X)$, put
\[
 a:=\gamma+u,
 \qquad
 f:=P_\gamma(u),
 \qquad
 \mu:=\MA_{\mathrm{poly},\mathbf R}(f).
\]
After a common finite refinement, $a$ is real-piecewise affine.  The
real-affine proof of the one-dimensional envelope statement in
\cite[Proposition~6.13]{BMPS26} gives
$f\in\mathrm{PA}_{\mathbf R}(\mathbf X)\cap\mathrm{PC}(\mathbf X)$
on that refinement.  Consequently $\mu$ is a nonnegative atomic measure
supported at its vertices.  A newly inserted vertex in the interior of an
old edge has two opposite primitive directions and equal inherited weights,
so its cone of nonnegative balancing relations is one-dimensional.  Thus the
hypothesis persists under the refinement.

We show that every point carrying positive $\mu$-mass lies in the contact
set $\{f=a\}$.  Let $v$ be a vertex of the refined structure, let $s_e$ be
the outgoing slope of $f$ along $e\in E(v)$, and suppose
\[
 m:=\mu\{v\}=\sum_{e\in E(v)}w_es_e>0,
 \qquad
 f(v)<a(v).
\]
For every incident edge choose a point $x_e$ at integral-affine distance
$\rho_e>0$ from $v$, so close that the open segment $(v,x_e)$ contains no
other vertex and $f$ is affine on $[v,x_e]$.  The continuous function
$a-f$ is positive at $v$, hence is bounded below by a positive constant on
a sufficiently small closed star; choose the points inside that star.

For $\delta>0$, define $f_\delta$ by raising the value at $v$ from $f(v)$
to $f(v)+\delta$, interpolating affinely from that new value to the unchanged
value $f(x_e)$ on every segment $[v,x_e]$, and leaving $f$ unchanged
elsewhere.  Choose $\delta$ smaller than the obstacle gap and so small that
\[
 0<\delta<
 \frac{m}{\sum_{e\in E(v)}w_e/\rho_e}.
\]
The new outgoing slope at $v$ along $e$ is
$s_e-\delta/\rho_e$.  Since $K_v=\mathbf R_{\geq0}w_v$, the local
polyhedral-convexity inequalities at $v$ reduce to the single inequality
\[
 \sum_e w_e\left(s_e-\frac{\delta}{\rho_e}\right)
 =m-\delta\sum_e\frac{w_e}{\rho_e}>0.
\]
At $x_e$, the slope on the segment from $v$ has decreased by
$\delta/\rho_e$ in the coordinate pointing away from $v$ and then returns
to the old slope.  Ordinary one-variable convexity is therefore preserved.
All other local convexity inequalities are unchanged.  The local
one-dimensional slope criterion, equivalently
\cite[Definitions~3.1--3.3 and Lemma~3.4]{BMPS26}, now shows that
$f_\delta$ is polyhedrally convex.

Moreover $f_\delta\leq a$, and $f_\delta-f$ has compact support, so
$f_\delta\leq\gamma+O(1)$.  Hence $f_\delta$ is an admissible competitor
for $P_\gamma(u)$ and satisfies $f_\delta(v)>f(v)$, contradicting the
definition of the envelope.  We conclude that
\[
 \operatorname{supp}\mu\subseteq\{f=a\}.
\]
The integrand $f-a$ vanishes on the support of the nonnegative atomic
measure $\mu$, and the asserted orthogonality integral is therefore zero.
The argument is local at the vertices and never removes a bounded edge or
cuts a graph cycle, proving the last two assertions.
\end{proof}

\begin{proof}[\Cref{cor:frontier-polyhedral-trivalent-graph}]
Balancing excludes a one-valent vertex.  At a two-valent vertex the two
primitive outgoing directions are opposite, and its positive balancing
relation spans $K_v$.  At a three-valent vertex, the three distinct outgoing
rays cannot all lie in a one-dimensional vector subspace.  Their span
therefore has dimension two, and the kernel of the direction map
$\mathbf R^3\to N_{\mathbf R}$ is one-dimensional.  Its positive balancing
vector spans $K_v$.  Thus every vertex satisfies $\dim K_v=1$, and
\Cref{thm:frontier-polyhedral-local-rank-one-graph} applies.
\end{proof}
